% Part: sets-functions-relations
% Chapter: functions
% Section: functions-relations

\documentclass[../../../include/open-logic-section]{subfiles}


\begin{document}

\olfileid[fa-IR]{sfr}{fun}{rel}

\olsection{توابع به‌منزلهٔ روابط}

\begin{explain}
تابعی که !!{element}s از~$A$ را به !!{element}s از~$B$
نگاشت می‌کند، آشکارا رابطه‌ای میان $A$ و~$B$ تعریف می‌کند؛ یعنی رابطه‌ای
که میان $x$ و $y$ برقرار است اگر و تنها اگر $f(x) = y$. در واقع، حتی ممکن است---
برای نمونه، اگر به تقلیل اجزای سازندهٔ ریاضیات علاقه‌مند باشیم---دست به
\emph{یکی‌انگاری} بزنیم و تابع~$f$ را با این رابطه، یعنی با مجموعه‌ای از
زوج‌ها، یکی بدانیم. آنگاه این پرسش پیش می‌آید:
کدام روابط به این شیوه توابع را تعریف می‌کنند؟
\end{explain}

\begin{defn}[گراف یک تابع] فرض کنید $f\colon A \to B$ یک تابع باشد.
\emph{گرافِ}~$f$ رابطهٔ $R_f \subseteq A \times B$ است که
به‌صورت زیر تعریف می‌شود:
\[
R_f = \Setabs{\tuple{x,y}}{f(x) = y}.
\]
\end{defn}

\begin{explain}
گراف یک تابع بنا بر اصل گسترش به‌طور یکتا تعیین می‌شود.
افزون بر این، اصل گسترش (برای مجموعه‌ها) بی‌درنگ
اصل ضمنی گسترش برای توابع را موجه می‌سازد؛
بنابر این اصل، اگر $f$ و~$g$ دامنه و هم‌دامنهٔ یکسانی داشته باشند، آنگاه اگر
برای همهٔ آرگومان‌ها مقادیر یکسانی داشته باشند، یکسان‌اند.

به همین ترتیب، اگر رابطه‌ای «تابع‌وار» باشد، آنگاه گراف یک تابع است.
\end{explain}

\begin{prop}\ollabel{prop:graph-function}
فرض کنید $R \subseteq A \times B$ چنان باشد که:
\begin{enumerate}
\item اگر $Rxy$ و $Rxz$، آنگاه $y = z$؛ و
\item برای هر $x \in A$ یک $y \in B$ وجود دارد چنان‌که $\tuple{x,
y} \in R$.
\end{enumerate}
آنگاه $R$ گراف تابع $f\colon A \to B$ است که به‌وسیلهٔ
$f(x) = y$ اگر و تنها اگر $Rxy$ تعریف می‌شود.
\end{prop}

\begin{proof}
فرض کنید یک $y$ وجود دارد چنان‌که $Rxy$. اگر $z \neq
y$ دیگری وجود داشت چنان‌که $Rxz$، شرط مربوط به~$R$ نقض می‌شد. پس اگر
یک $y$ وجود داشته باشد چنان‌که $Rxy$، این $y$ یکتا است و بنابراین $f$
خوش‌تعریف است. آشکارا $R_f = R$.
\end{proof}

\begin{explain}
هر تابع $f\colon A \to B$ گرافی دارد؛ یعنی رابطه‌ای روی $A
\times B$ که با $f(x) = y$ تعریف می‌شود. از سوی دیگر، هر رابطهٔ~$R
\subseteq A \times B$ که ویژگی‌های داده‌شده در
\olref{prop:graph-function} را داشته باشد، گراف تابعی~$f \colon A \to
B$ است. به‌سبب این پیوند نزدیک میان توابع و
گراف‌هایشان، می‌توانیم یک تابع را صرفاً گراف آن در نظر بگیریم. به
بیان دیگر، توابع را می‌توان با روابطی معین، یعنی با
مجموعه‌هایی معین از چندتایی‌ها، یکی انگاشت. \oliflabeldef{sfr:rel:ref:sec}{بااین‌حال،
توجه کنید که روح این «یکی‌انگاری» همان است که در
\olref[sfr][rel][ref]{sec} آمده است: این ادعایی دربارهٔ متافیزیک
توابع نیست، بلکه مشاهده‌ای است مبنی بر اینکه مناسب است توابع را همچون
مجموعه‌هایی معین \emph{در نظر بگیریم}. یکی از دلایل این مناسب‌بودن آن است
که اکنون می‌توانیم}{اکنون می‌توانیم} عملیات مشابهی را روی
توابع انجام دهیم که روی روابط انجام دادیم (نگاه کنید به
\olref[sfr][rel][ops]{sec}). به‌ویژه:
\end{explain}

\begin{defn}\ollabel{defn:funimage}
فرض کنید $f \colon A \to B$ تابعی باشد و $C\subseteq A$ برقرار باشد.

\emph{تحدیدِ}~$f$ به~$C$ تابع
~$\funrestrictionto{f}{C}\colon C \to B$ است که به‌صورت زیر تعریف می‌شود:
$(\funrestrictionto{f}{C})(x) = f(x)$ برای هر $x \in C$. به
بیان دیگر، $\funrestrictionto{f}{C} = \Setabs{\tuple{x, y} \in R_f}{x \in
C}$.

\emph{اعمالِ}~$f$ بر~$C$ برابر است با $\funimage{f}{C} =
\Setabs{f(x)}{x \in C}$. این را \emph{تصویرِ}~$C$
تحت~$f$ نیز می‌نامیم.
\end{defn}

\begin{explain}
از این تعریف‌ها نتیجه می‌شود که $\ran{f} =
\funimage{f}{\dom{f}}$ برای هر تابع~$f$ برقرار است.
\oliflabeldef{sfr:rel:ops:sec}{این مفاهیم با توجه به تعریف‌های
\olref[sfr][rel][ops]{sec} و یکی‌انگاری توابع با روابط، دقیقاً
همان‌اند که انتظار می‌رود. اما دو عمل دیگر---وارون‌ها و حاصل‌ضرب‌های نسبی---به
جزئیات اندکی بیشتر نیاز دارند. آن جزئیات را در
\olref[inv]{sec} و
\olref[cmp]{sec} ارائه خواهیم کرد.}{}
\end{explain}

\end{document}
